% This file should be replaced with your file with thesis content.
%=========================================================================
% Authors: Michal Bidlo, Bohuslav Křena, Jaroslav Dytrych, Petr Veigend and Adam Herout 2019

% For compilation piecewise (see projekt.tex), it is necessary to uncomment it and change
% \documentclass[../projekt.tex]{subfiles}
% \begin{document}

% ===================================
%           INTRO
% ===================================

\chapter{Introduction}
Classification of speech and music involves determining whether a signal, whether whole or just a segment, is speech or music.
The focus of this work is on-line classification, the process of classifying a continuous live stream of data~\cite{Eyben2016}.
In contrast to the offline approach, data are processed while being recorded, and the result is returned immediately after the computation is complete.

This task is typically a crucial preprocessing step in larger pipelines, such as automatic speech recognition~\cite{BENZEGHIBA2007763} or telephony; therefore, the aim is not only to improve the precision of the classifier but also to enhance its efficiency.
Historically, complexity constraints have led researchers to develop relatively simple machine learning models, such as Decision Trees~\cite {Lavner2009-zo}, Gaussian Mixture Models (GMMs), or Support Vector Machines (SVMs)~\cite{KHONGLAH201671}.

The improvement in computational power of devices that utilize applications requiring speech/music classifiers, combined with the rise of deep learning, has led to a resurgence and advancement in the field of speech and music classification in recent years.
Researchers have experimented with various neural network architectures, primarily Convolutional Recurrent Neural Networks (CRNNs)~\cite{choi2016convolutionalrecurrentneuralnetworks}.

Many existing solutions for on-line speech and music classification are either proprietary or limited in scope. 
Common limitations include restrictive assumptions about the input signal, such as the presence of only speech or only music, and a lack of training on more complex acoustic scenarios, for example overlapping speakers or mixed audio content. 
Other approaches are overly general and therefore inefficient or impractical when deployed specifically for this task.

The objective of this thesis is to review existing machine learning and deep learning methods and to propose a novel deep learning–based approach for on-line speech and music classification.

The experiments have shown ... 

The thesis is organized as follows.
\Cref{overview} describes the historical background, defines the concept of on-line speech and music classification, and presents the constraints on the methods used, as well as the motivation for this thesis.
\Cref{cap:classification_approaches} provides an overview of existing methods, their architectures and limitations, and presents the models used as references for this work, including their strengths, weaknesses, and the adjustments made to the original implementations.
\Cref{cap:dataset} covers the dataset used in the training and evaluation of the classifiers.

% ===================================
%           THEORY
% ===================================

\chapter{Background and Problem Context}
\label{overview}
This chapter describes how the task of speech and music classification was first introduced to the public and how it evolved into the building block of most applications, followed by the definition and constraints of the on-line classifier.

\section{Historical Background}
Automatic classification of speech and music originated as a practical quality-of-life feature in radio broadcasting, enabling listeners to mute advertisements or switch between channels automatically~\cite {543290}. 

The surge of ASR applications has made the audio classification task boom, as a need to detect speech in environments with much noise or music, because the applications would run unnecessarily or incorrectly~\cite {596192}

Today, speech–music classification is a fundamental component of modern audio processing systems. 
It is used in voice-activated assistants, music streaming platforms, and telecommunication systems, where accurate content identification enables appropriate codec selection and efficient resource usage~\cite{ren2024fewertokenneuralspeechcodec, jia2025practicalrealtimeneuralvideo}. 
Beyond consumer applications, it also supports large-scale infrastructures such as broadcast monitoring and audio archiving.

\section{On-Line Classification}
\label{sec:on-line_classification}
This thesis will use the definitions of the term 'on-line processing' according to Eyben's doctoral thesis~\cite{Eyben2016}.
On-line processing implicitly requires real-time processing; therefore, both terms are introduced.

\textit{"Real-time processing refers to processing of data in real-time. This implies two things: (a) the processing of the data requires less time than the duration of the data, i.e., the time it will take to record or play-back the data at normal speed, and (b) the processing lag is as low as possible, ideally zero."}

\textit{"On-line processing refers to processing of a continuous live stream of data while data are recorded or transmitted, i.e., computing and returning results continuously before the end of the data stream has been seen."}

After these definitions, on-line classification can be described as a process in which incoming data are processed continuously and classified into predefined classes without access to future samples, while maintaining low processing lag.

In practice, on-line processing usually requires the use of a buffer to aggregate incoming samples.
Once the buffer is filled, the classification is performed on the aggregated segment.
This aggregation introduces a delay between data acquisition and the availability of classification results, which is directly determined by the buffer size.
Consequently, the requirement for low latency limits the buffer duration to short time intervals, typically in the range of 10-30 ms.
For the use of features that require a longer temporal context, such as statistics computed over multiple frames, an additional buffer containing previous samples may be employed; however, the processing is still triggered by the primary buffer, and the resulting lag is therefore not eliminated. 

Buffering is not the only source of delay in an on-line classification system.
Additional latency is introduced by the computational cost of feature extraction and classification, which must be completed after each buffer is processed.
Therefore, high computational complexity can significantly increase overall processing lag and compromise real-time operation.
For this reason, on-line classification requires the use of computationally efficient algorithms with predictable execution times.

\subsection{Constraints}
\label{subsec:constraints}
Based on the above considerations, the following constraints apply to the on-line classification approach used in this work:
\begin{itemize}
    \item \textbf{Causality}: The classification must be performed using only past and present samples, without access to future data.
    \item \textbf{Real-time operation}: All processing steps must be completed within the time duration of the processed data segment.
    \item \textbf{Low processing lag}: The delay between data acquisition and the availability of classification results must be minimized.
    \item \textbf{Limited buffer size}: Buffer durations must be kept short to satisfy latency requirements.
    \item \textbf{Bounded computational complexity}: Feature extraction and classification must have predictable and sufficiently low computational cost.
\end{itemize}

%\section{Motivation}
%\textcolor{green}{?}

%This thesis is motivated by the lack of modern open-source speech and music classification models designed specifically for on-line speech and music classification.
%Many existing methods are either proprietary or intended for more general audio classification tasks, which often leads to unnecessarily high model complexity when applied to this problem.


\chapter{Classification Approaches}
\label{cap:classification_approaches}
The following chapter describes the general architecture of existing classifiers, presents models selected as references to the proposed solution, and adjustments made to these models to make them fit the task of on-line speech and music classification.

\section{Overview}

\begin{figure}[htbp]
\includegraphics[width=\textwidth]{docs/figures/general_classifier.pdf}
\label{fig:general_classifier}
\caption{Example diagram of generic speech/music classifier}
\end{figure}

Classifiers are algorithms whose goal is to assign an input sample to one of a set of predefined classes. 

Figure~\ref{fig:general_classifier} illustrates a minimal processing pipeline of a generic classification system.

The first stage of the pipeline consists of pre-processing steps, which are largely influenced by the architecture of specific classifier and its application. 
In the case of on-line audio classification, buffering and segmentation of the input signal are required due to the constraints described in \cref{subsec:constraints}. 

Raw input samples do not provide a robust representation of an audio signal, which motivates the use of feature extraction in the majority of audio classification pipelines.

Although a speech and music classifier can operate as a standalone application, audio classifiers are most commonly used as pre-processing components in larger processing pipelines. 
As a module in this context, classifiers often serve to prevent subsequent processing stages from wasting computational resources or from operating on inactive or invalid input.

\section{Traditional Methods}
Traditional methods of speech and music classification refer to usage of machine learning algorithms such as Gaussian Mixture Models, Decision Trees or Support Vector Machines.

The advantage of using machine learning models in comparison to deep learning algorithms is the low computational complexity in both training and inference stages. 

Main drawback of traditional methods is the need for selection of right features to represent input signal as they cannot work on raw input.
Selection is carried out by thorough experiments with combinations of features or integration of the selection algorithm, such as sequential forward/backward floating selection~\cite{PUDIL19941119}, for the selection of best features by eliminating features with low variance.

\section{Reference Machine Learning Models}
Models in this subsection were picked as references to the proposed method.

\subsection{Decision Tree Classifier by Lavner and Ruinskiy~\cite{Lavner2009-zo}}
This method proposed a decision-tree-based algorithm for speech and music classification.

Feature extraction and selection for this model is implemented as described by the paper.
Authors assume every input segment is either speech or music; silent and noisy frames must therefore be eliminated from the training phase of the model.
For the fairness of comparison, all models were evaluated on 2-class and 3-class classification task.

\textcolor{green}{separate tables?}

\subsection{Gaussian Mixture Model and Support Vector Machine by Khonglah and Prasanna~\cite{KHONGLAH201671}}
\textcolor{green}{!}
\section{Deep Learning Methods}

\section{Reference Deep Learning Models}

\subsection{YamNET~\cite{drossos2020soundeventdetectiondepthwise}}


%Most classifiers do not operate directly on raw input samples, as these typically do not provide a compact or robust representation of the signal characteristics. 
%Instead, a feature extraction stage is employed to transform the buffered audio segments into a set of descriptive features that capture relevant temporal and spectral properties of the signal. 
%Commonly used features in speech and music classification include spectral features, energy-based measures, and statistics computed over short time windows.

%The extracted features are then passed to the classification model, which maps the feature representation to one of the predefined classes. 
%Various model architectures can be used at this stage, ranging from traditional machine learning approaches to modern neural network-based methods. 
%The choice of the classification model is influenced by factors such as classification accuracy, computational complexity, and suitability for on-line operation.

%Finally, the classifier output may be post-processed to obtain more stable decisions over time. 
%This step can include temporal smoothing or aggregation of multiple frame-level decisions into a single segment-level output. 
%In speech and music classification, the final result is typically a binary decision indicating whether the input signal corresponds to speech or music.





% ===================================
%           DATASET
% ===================================

\chapter{Dataset}
\label{cap:dataset}

This chapter describes the dataset used in this work, the requirements imposed on datasets for speech and music classification, the construction of a custom dataset and its properties.

\section{Overview}

The dataset used in training and evaluation of classifier should represent real-world scenarios as accurately as possible. 
For this reason, it is useful to define the range of acoustic events that the classifier may encounter in practice.

Al-Shoshan~\cite{ALSHOSHAN200695} proposes the following categorization of audio signals:
\begin{enumerate}
    \item \textit{"Speech signal composed of a single talker within a specific time period."}
    \item \textit{"Pure music signal without any speech component."}
    \item \textit{"Mixture of single-talker speech and background music."}
    \item \textit{"Songs, i.e., a mixture of music and a singer's voice."}
    \item \textit{"Singing without musical accompaniment."}
    \item \textit{"Abnormal music, including unconventional cadences, human whistles, or non-musical sounds embedded as part of a musical structure."}
    \item \textit{"Speech signal composed of two or more speakers talking simultaneously."}
    \item \textit{"Non-speech and non-music signals, such as car, motor, or fan sounds."}
    \item \textit{"Complex sound mixtures involving multiple speakers, singers, and musical sources."}
\end{enumerate}

This categorization covers both natural and artificially generated sound types that may occur in real-world audio signals and therefore provides a useful reference for defining the scope of the dataset used in this work.

\section{Composition and Construction}

Following this categorization, three main classes of data can be defined:
\begin{enumerate}
    \item \textbf{Speech}: Categories 1, 3, and 7.
    \item \textbf{Music}: Categories 2, 4, 5, 6, and 9.
    \item \textbf{Non-speech Non-music audio}: Category 8 and silence, this class will be labeled as \textbf{Inactive}. 
\end{enumerate}


%For the purposes of this work, these categories are mapped to three output classes of the classifier as follows:
%\begin{itemize}
%    \item Categories 1, 3, and 7 correspond to the \textit{speech} class.
%    \item Categories 2, 4, 5, 6, and 9 correspond to the \textit{music} class.
%    \item Category 8, together with silence segments, corresponds to the \textit{inactive} class.
%\end{itemize}

The dataset used in the training and evaluation phases for all classifiers mentioned in this thesis is a combination of publicly available datasets, mainly from huggingface.com.

Sources were chosen to take into account all categories of audio signals described in the previous section and will now be enumerated.

\begin{enumerate}
    \item People's speech (MLCommons) 
    
One of the largest public datasets used for English speech recognition, hosted on huggingface.com.
It contains 30,000 hours of clean speech divided into recordings around 15 seconds.
It was chosen as a majority part of clean speech part of the final dataset.
16,000 rows were taken from this dataset, corresponding to approximately 67 hours of audio.

    \item English--Arabic Speech Translation (Ammagra)

This dataset was chosen to provide the rest of the clean speech.
It is hosted on huggingface.com and contains over 300,000 rows of recordings mostly under 10 seconds.
4,000 rows, which correspond to roughly under 8 hours.

    \item 

    
\end{enumerate}










The dataset used in this thesis is a combination of publicly available datasets listed below:


Weak labels were converted into strong, time-aligned annotations by automatically detecting active and silent regions. Class-specific activity detection was applied:

\begin{itemize}
    \item \textbf{Speech:} Silero VAD~\footnote{\url{https://github.com/snakers4/silero-vad}} for the detection of speech activity.
    \item \textbf{Music:} PyDub~\footnote{\url{https://github.com/jiaaro/pydub}} amplitude-based silence detection.
\end{itemize}

\section{Statistics}
\label{sec:dataset_statistics}
All data were randomly shuffled and split into disjoint subsets:

\begin{itemize}
    \item \textbf{Training:} 80\%
    \item \textbf{Validation:} 10\%
    \item \textbf{Test:} 10\%
\end{itemize}

\begin{table}[ht]
    \centering
    \caption{Dataset Statistics and Composition}
    \label{tab:dataset_stats}
    \begin{tabular}{lrrr}
        \toprule
        \textbf{Class / Subclass} & \textbf{Frames} & \textbf{Duration (min)} & \textbf{Share (\%)} \\
        \midrule
        \textbf{Speech (Total)} & \textbf{20,187,451} & \textbf{3,364.58} & \textbf{65.8\%} \\
        \hspace{3mm} Clean & 19,769,776 & 3,294.96 & 64.4\% \\
        \hspace{3mm} Noisy & 284,171 & 47.36 & 0.9\% \\
        \hspace{3mm} Noisy Environment & 133,504 & 22.25 & 0.4\% \\
        \addlinespace
        \textbf{Music (Total)} & \textbf{10,074,118} & \textbf{1,679.02} & \textbf{32.8\%} \\
        \hspace{3mm} Mainstream & 5,436,783 & 906.13 & 17.7\% \\
        \hspace{3mm} Acoustic & 2,856,838 & 476.14 & 9.3\% \\
        \hspace{3mm} Instrumental & 1,187,275 & 197.88 & 3.9\% \\
        \hspace{3mm} Vocal & 470,332 & 78.39 & 1.5\% \\
        \hspace{3mm} Electronic & 122,890 & 20.48 & 0.4\% \\
        \addlinespace
        \textbf{Noise (Total)} & \textbf{439,630} & \textbf{73.27} & \textbf{1.4\%} \\
        \midrule
        \textbf{Total} & \textbf{30,701,199} & \textbf{5,116.87} & \textbf{100.0\%} \\
        \bottomrule
    \end{tabular}
\end{table}



% ===================================
%           PIPELINE
% ===================================

%\chapter{Pipeline}
%\label{cap:pipeline}

% ===================================
%           MODELS
% ===================================
%\chapter{Models}
%\label{cap:ref_models}

% ===================================
%           EXPERIMENTS
% ===================================
%\chapter{Experiments}
%\label{cap:experiments}

% ===================================
%           CONCLUSION
% ===================================
%\chapter{Conclusion}
%\label{cap:conclusion}


%=========================================================================

% For compilation piecewise (see projekt.tex), it is necessary to uncomment it
% \end{document}
